\chapter{Engineering a Real-Time Estimator in C++20}
\label{ch:engineering}
\chaptermark{Engineering in C++20} % short running header
\sloppy % long code identifiers cannot be hyphenated

\begin{objectives}
\begin{itemize}
  \item Know the C++ facts an estimator depends on: stack and heap, object lifetime, and the C++20 features the project uses.
  \item Understand why a real-time estimator must be deterministic, and how the project's memory tiers achieve it.
  \item Know how Eigen represents matrices, blocks and expressions, and use fixed-max-size matrices, member workspaces and
        allocation guards correctly.
  \item Avoid the Eigen pitfalls that silently corrupt numbers: dangling expression templates and aliasing.
  \item Understand what the sanitizers check and how, and how a target-based CMake build is organized.
  \item Design a test strategy in which every test can fail, and a systematic procedure for debugging a diverging filter.
\end{itemize}
\end{objectives}

The mathematics of the previous chapters is necessary but not sufficient. An estimator that is mathematically correct can still
miss its deadline because of an allocation, produce garbage because of an Eigen expression template, or pass all its tests because
the tests cannot fail. This chapter collects the engineering practices that the project encodes in AGENTS.md.

% ------------------------------------------------------------------------------
\section{C++ essentials for estimator code}
\label{sec:cpp-essentials}

\subsection{Stack, heap and object lifetime}

A C++ program keeps its data in two kinds of memory. The \term{stack} holds the local variables of the functions currently running. Entering a
function reserves a block of it by moving a pointer, and leaving the function releases the block by moving the pointer back: the cost is a
few instructions, it is the same every time, and the size of each block is fixed at compile time. The \term{heap} holds objects whose size or
lifetime is decided at run time. Obtaining heap memory (\code{new}, \code{std::vector} growth, a dynamic-size Eigen matrix) calls the memory
allocator, which must find a free block of the right size, may take a lock shared with other threads, and occasionally asks the operating
system for more pages. Its cost is usually modest and occasionally large, which is exactly what a real-time loop cannot afford.

Every C++ object has a lifetime that ends at a precisely defined point: a local variable when its scope closes, a member when its owner is
destroyed. At that point its destructor runs. Tying resources (memory, files, locks) to object lifetime is called \term{RAII} (resource
acquisition is initialization): a \code{std::vector} frees its heap block in its destructor, so it cannot leak. The project uses RAII
everywhere and never calls \code{delete} by hand. The stack is small (typically \SI{1}{MB} on Windows, \SI{8}{MB} on Linux), so large
matrices belong in class members created once, not in function locals.

\subsection{The C++20 features the project uses}

\begin{lstlisting}
inline constexpr int kImuDim = 15;   // compile-time constant, one definition

[[nodiscard]] bool update(const PoseMeasurement& z) noexcept; // result must be used
void ingest(std::span<const ImuSample> samples) noexcept;     // view of contiguous data
\end{lstlisting}

\begin{itemize}
  \item \code{constexpr} values are computed at compile time and can size arrays and Eigen matrices; \code{inline constexpr} in a header gives one
        shared definition. Dimensions, indices (AGENTS.md C3) and the $\chi^2$ table are \code{constexpr}.
  \item \code{[[nodiscard]]} makes the compiler warn (with \code{-Werror}, fail) when a caller ignores the return value: a rejected update or a failed
        factorization must not be silently dropped.
  \item \code{noexcept} promises that a function throws no exception. If one escapes anyway the program terminates, so the promise is only made
        on the hot path, where errors are reported through return values instead.
  \item \code{std::span<const T>} is a non-owning view (a pointer and a length) of contiguous elements. A function taking a span accepts an array, a
        \code{std::vector} or part of one without copying and without caring who owns the memory.
  \item \code{std::variant<A, B>} holds exactly one of several types, without heap allocation. \code{std::visit} calls the matching overload:
\end{itemize}

\begin{lstlisting}
template <class... Ts> struct Overloaded : Ts... { using Ts::operator()...; };
template <class... Ts> Overloaded(Ts...) -> Overloaded<Ts...>;

for (const SensorEvent& event : events) {           // one presorted stream
  std::visit(Overloaded{
      [&](const ImuSample& s)  { filter.propagate(s); },
      [&](const ImageEvent& e) { pipeline.on_image(e); },
  }, event);
}
\end{lstlisting}

\code{Overloaded} inherits the call operators of all the lambdas passed to it, so \code{std::visit} finds the one whose parameter type matches the
alternative currently stored. Adding a new event type to the variant without handling it is a compile error, not a silent omission.

% ------------------------------------------------------------------------------
\section{Determinism and the memory tiers}

A flight controller consumes the estimator's output at a fixed rate. A late estimate is a wrong estimate. The main source of
unpredictable latency in C++ is dynamic memory allocation: \code{new}, a \code{std::vector} that grows, a dynamic-size Eigen temporary. The
allocator may take locks, search free lists or ask the operating system for pages, with latency that is usually small and
occasionally large. The project therefore separates three tiers (AGENTS.md E4).

\begin{center}\small
\begin{tabularx}{\textwidth}{@{}l X X@{}}
\toprule
Tier & Scope & Rule \\
\midrule
Load & parsing, building event arrays, constructing the filter & allocate freely, before streaming starts \\
Core & \code{vio\_core}, per IMU sample and per update & zero heap allocations \\
Frontend & OpenCV tracking & may allocate, must reuse buffers; allocations are measured and logged \\
\bottomrule
\end{tabularx}
\end{center}

\subsection{Fixed-max-size matrices}

The MSCKF covariance changes size as clones come and go, but never exceeds $15 + 6N_{max}$. Eigen supports exactly this case: a
matrix whose \emph{runtime} size is dynamic but whose \emph{maximum} size is a compile-time constant. Its storage is an inline array of the
maximum size, so resizing within the bound never allocates.

\begin{lstlisting}
inline constexpr int kMaxClones = 10;                  // decided in M10a
inline constexpr int kMaxDim    = 15 + 6 * kMaxClones;

using CovMatrix = Eigen::Matrix<double, Eigen::Dynamic, Eigen::Dynamic,
                                Eigen::ColMajor, kMaxDim, kMaxDim>;

CovMatrix P;          // inline storage for 75 x 75 doubles, no heap
P.resize(15, 15);     // allowed: within the compile-time maximum
\end{lstlisting}

Large workspaces ($S$, $K$, stacked $H$, temporaries for $\Phi P$) are class members, created once, rather than function-local
variables. Function-local fixed-max matrices would not allocate either, but they would live on the stack, and several of them at
once can exceed it (\cref{ch:cloning}).

\subsection{Proving it: allocation guards}

``It shouldn't allocate'' is not evidence. The project overrides the global \code{operator new} in its tests to count allocations
(\code{tests/support/alloc\_guard.hpp}), and wraps every core hot-path test with a guard that asserts a count of zero. In debug builds,
\code{EIGEN\_RUNTIME\_NO\_MALLOC} together with \code{Eigen::internal::set\_is\_malloc\_allowed(false)} makes Eigen itself abort on any internal
allocation and reports where it happened. Both mechanisms are needed, because they see different allocations. The guard sees everything that goes
through \code{operator new}: \code{std::vector} growth, \code{std::string}, \code{std::function}, \code{std::make\_unique}. Eigen, however, obtains its heap memory
directly from \code{std::malloc} through its own aligned allocator, so a dynamic \code{Eigen::MatrixXd} is invisible to the guard. Only Eigen's flag catches it,
and only in the \code{debug-sanitize} preset, where the flag is defined.

% ------------------------------------------------------------------------------
\section{Eigen in brief}
\label{sec:eigen-basics}
\index{Eigen}\index{expression templates}

\paragraph{The matrix type.} Every dense Eigen matrix is an instance of one class template,
\begin{lstlisting}
Eigen::Matrix<Scalar, Rows, Cols, Options, MaxRows, MaxCols>
\end{lstlisting}
\code{Rows} and \code{Cols} are either compile-time numbers (a \emph{fixed-size} matrix such as \code{Eigen::Matrix3d}, stored inline with no heap) or
\code{Eigen::Dynamic} (size chosen at run time, normally on the heap, such as \code{Eigen::MatrixXd}). \code{MaxRows} and \code{MaxCols} give a dynamic
matrix a compile-time upper bound, which turns its storage into an inline array: the fixed-max-size matrices used below. \code{Options} chooses the
storage order.

\paragraph{Storage order.} Memory is one-dimensional. A \emph{column-major} matrix (Eigen's default) stores column after column, so entry $(i, j)$ of
an $m\times n$ matrix lives at offset $i + jm$; a \emph{row-major} matrix stores row after row, offset $j + in$. Loops are fastest when they walk
memory contiguously, so with column-major storage the inner loop should run down a column. Mixing orders is legal but may force copies.

\paragraph{Blocks.} \code{P.block(i, j, r, c)}, \code{P.topLeftCorner(15, 15)}, \code{v.segment<3>(6)} and their relatives return \emph{views}: they refer to
the original storage, copy nothing, and can be assigned to. \code{P.block<3, 3>(6, 6) = ...} writes the attitude block of $P$ in place. The fixed-size
forms (template arguments) let Eigen unroll small operations completely.

\paragraph{Passing matrices and blocks to functions.} A parameter of type \code{const Eigen::Matrix3d\&} accepts only that exact type. To accept any
matrix or block of compatible layout, use \code{Eigen::Ref<const Eigen::MatrixXd>} (read-only) or \code{Eigen::Ref<Eigen::MatrixXd>} (writable). A
read-only \code{Ref} quietly makes a temporary copy when the argument's layout is incompatible, and for a dynamic type that copy allocates; a writable
\code{Ref} refuses to compile instead. The quiet copy is an Eigen allocation, so \code{EIGEN\_RUNTIME\_NO\_MALLOC} catches it and the
\code{operator new} guard does not.

\paragraph{Wrapping existing memory.} \code{Eigen::Map<const Eigen::Vector3d> v(ptr)} views three doubles at \code{ptr} as a vector, without copying.
It is the right tool for parsed buffers, and the wrong one for quaternions read from EuRoC: \code{Eigen::Map<Eigen::Quaterniond>} interprets memory as
$(x, y, z, w)$, while the CSV stores $(w, x, y, z)$ (AGENTS.md C2).

\paragraph{Expression templates.} \code{A + B} does not compute a matrix. It returns a small object recording ``the sum of \code{A} and \code{B}''. Only when
that object is assigned to a matrix does Eigen run a loop, and it fuses the whole expression into it: \code{D = A + B + C} is one loop with no
temporaries. Products are the exception: they are evaluated into a temporary unless told otherwise (below). This design makes Eigen fast, and it
is also the source of the \code{auto} pitfall described next.

\paragraph{Decompositions.} Linear systems are solved through decomposition objects: \code{S.llt().solve(B)} (Cholesky, for symmetric positive definite
$S$), \code{.ldlt()} (robust for semi-definite), \code{.householderQr()} and \code{.colPivHouseholderQr()} (QR, \cref{ch:numlinalg}). A decomposition object owns
a copy of the factor; for fixed-max matrices that copy is inline. \code{Eigen::LLT<Eigen::Ref<Eigen::MatrixXd>{}> llt(A)} factors \code{A} \emph{in place},
overwriting it with the factor and allocating nothing. Symmetric structure can be exploited with \code{P.selfadjointView<Eigen::Lower>()}, which reads
only one triangle.

% ------------------------------------------------------------------------------
\section{Eigen pitfalls}

\paragraph{Expression templates and \code{auto}.} Eigen operations return lightweight \emph{expression} objects that are evaluated
lazily, when they are assigned to a matrix. An expression holds the matrices it uses by reference and nested expressions by value. Storing one
with \code{auto} therefore has two dangers. If a referenced matrix is a temporary (made by \code{.eval()}, or returned by value from a function),
the reference dangles. And even if every operand still exists, the computation happens only at the later assignment, with whatever values
the operands have then.

\begin{lstlisting}
auto bad = (a + b).eval().transpose(); // refers to the temporary made by eval()
Eigen::RowVector3d x = bad;            // undefined behavior: that temporary is gone
Eigen::RowVector3d good = (a + b).transpose(); // concrete type: evaluated now

auto lazy = A * v;                     // nothing is computed here
A(0, 0) = 10.0;
Eigen::Vector3d y = lazy;              // computed now, with the modified A
\end{lstlisting}

The rule is simple: never use \code{auto} with Eigen expressions (AGENTS.md E5). Name the concrete type, or call \code{.eval()} when a
temporary really is wanted. AddressSanitizer reports the dangling case (as a use after scope), but not the stale one, which reads only valid
memory (Exercise~30.2).

\paragraph{Aliasing.} Eigen assumes that the destination of a matrix product may alias its operands and evaluates products into a
temporary. With dynamic sizes that temporary allocates. When the destination is known not to alias, \code{.noalias()} writes directly
into it. When it does alias, the temporary is required: write into a member workspace explicitly.

\begin{lstlisting}
tmp_15xN_.noalias() = Phi * P.topRightCorner(15, nc);  // reads P, writes workspace
P.topRightCorner(15, nc) = tmp_15xN_;                   // copy back
P.bottomLeftCorner(nc, 15) = tmp_15xN_.transpose();     // keep symmetry
\end{lstlisting}

Coefficient-wise operations do not alias-check. \code{A = A.transpose()} is the classic bug; use \code{A.transposeInPlace()} or symmetrize
with a workspace.

\paragraph{Passing and alignment.} Pass Eigen objects by \code{const\&} (or \code{Eigen::Ref} for blocks). Do not add
\code{EIGEN\_MAKE\_ALIGNED\_OPERATOR\_NEW}: since C++17, aligned \code{new} handles over-aligned types. The real alignment hazard is mixing
\code{-march} or alignment-changing flags between translation units or with dependencies, which the project forbids.

\paragraph{Decompositions without allocation.} Eigen's decompositions normally own their storage. Fixed-max decomposition objects held as
members, or the in-place variants that factor a matrix in its own storage, avoid per-call allocation. Whatever the API, verify with
\code{EIGEN\_RUNTIME\_NO\_MALLOC}, which sees Eigen's allocations where the \code{operator new} guard cannot.

% ------------------------------------------------------------------------------
\section{Time and events}

\begin{lstlisting}
using Clock = std::chrono::nanoseconds;
struct ImuSample   { Clock t; Eigen::Vector3d gyro, accel; };
struct ImageEvent  { Clock t; int index; };
using SensorEvent = std::variant<ImuSample, ImageEvent /*, PoseMeasurement */>;

// Delta t: subtract in integers, convert only the difference.
const double dt = std::chrono::duration<double>(t1 - t0).count();
\end{lstlisting}

Absolute time stays in 64-bit integer nanoseconds (\cref{sec:floating}). Events form one presorted sequence built at load time, so the
streaming loop is a single pass over contiguous memory, dispatched with \code{std::visit} as in \cref{sec:cpp-essentials}: no priority queue, no
allocation, no reordering surprises.

% ------------------------------------------------------------------------------
\section{Numerical robustness checklist}

\begin{itemize}
  \item Normalize every quaternion after every multiplicative update (propagation, injection, cloning).
  \item Symmetrize $P$ after propagation and after every update.
  \item Use the Joseph form; factor $S$ with Cholesky or LDL$^\mathsf{T}$; treat a failed factorization as an error to report, not to hide.
  \item Guard every $0/0$ with Taylor branches; use $\atantwo$ for angles; handle $\Log$ near $\pi$.
  \item Reject features with non-positive or tiny depth, poor parallax or large reprojection error before linearizing.
  \item Count and log gate rejections and triangulation failures.
  \item Never clamp negative eigenvalues of $P$ to hide a problem; find out why they appeared.
  \item In debug builds, assert finiteness at module boundaries so that a \code{NaN} is caught where it is born.
\end{itemize}

% ------------------------------------------------------------------------------
\section{Sanitizers}
\label{sec:sanitizers}
\index{sanitizers}\index{AddressSanitizer}\index{UndefinedBehaviorSanitizer}

C++ leaves many errors undefined: reading past the end of an array, using memory after freeing it, overflowing a signed integer. Undefined
behavior does not mean a crash; it means the compiler may assume it never happens, so the program can produce plausible wrong numbers, behave
differently between builds, or fail far from the bug. Sanitizers are compiler-inserted checks that turn these silent errors into immediate,
precise reports.
\begin{itemize}
  \item \textbf{AddressSanitizer} (ASan, \code{-fsanitize=address}) keeps a \emph{shadow memory}, one byte for every eight bytes of program memory,
        recording which bytes may be accessed. Every heap block, stack variable and global is surrounded by poisoned \emph{redzones}, and freed
        memory stays poisoned for a while. The compiler inserts a shadow check before every load and store. ASan reports heap, stack and
        global buffer overflows, use after free, use after scope and double frees, with stack traces of the access and of the allocation. The
        cost is roughly a factor of two in speed and three in memory.
  \item \textbf{UndefinedBehaviorSanitizer} (UBSan, \code{-fsanitize=undefined}) inserts checks for signed integer overflow, out-of-range shifts, integer
        division by zero, null and misaligned pointer use, invalid \code{bool} and \code{enum} values, and floating-point-to-integer conversions that
        overflow.
\end{itemize}
By default UBSan prints a message and \emph{continues}, so a test that triggers it still passes. \code{-fno-sanitize-recover=all} makes every report
abort the process, and the test preset sets \code{UBSAN\_OPTIONS=print\_stacktrace=1:halt\_on\_error=1} so the report carries a stack trace
(AGENTS.md E3). Frame pointers (\code{-fno-omit-frame-pointer}) make those traces reliable. On Windows, LeakSanitizer does not exist, so leaks are
not detected (RAII makes them rare), and an instrumented executable needs the ASan runtime DLL from \code{C:\textbackslash msys64\textbackslash clang64\textbackslash bin}
on its \code{PATH}. Sanitizers are a debug-build tool: the \code{release} preset builds without them, and both presets must pass.

% ------------------------------------------------------------------------------
\section{Modern CMake in brief}
\label{sec:cmake}
\index{CMake}

CMake describes a build as a set of \emph{targets} (libraries and executables) with properties. Modern CMake attaches every setting to a target
rather than to global variables, and lets settings propagate along the dependency graph through three keywords:
\code{PRIVATE} (used to build this target only), \code{INTERFACE} (used only by targets that link to this one) and \code{PUBLIC} (both). A library's public
include directory is \code{PUBLIC}, because its users include its headers; a warning flag is \code{PRIVATE}, because it concerns only the library's
own code.

\begin{lstlisting}
add_library(vio_warnings INTERFACE)            # carries flags, builds nothing
target_compile_options(vio_warnings INTERFACE -Wall -Wextra -Wconversion -Werror)

add_library(vio_core STATIC src/so3.cpp src/es_ekf.cpp)
target_include_directories(vio_core PUBLIC include)
target_link_libraries(vio_core PUBLIC Eigen3::Eigen PRIVATE vio_warnings)
\end{lstlisting}

The \code{vio\_warnings} target is how warnings stay off the dependencies (AGENTS.md E3): only project targets link it. Dependencies fetched with
\code{FetchContent} at a pinned tag are declared \code{SYSTEM}, so their headers are included with \code{-isystem} and do not trigger the project's warnings.
\code{CMakePresets.json} stores named configurations (generator, compiler, build type, flags, test environment), so that
\code{cmake --preset debug-sanitize}, \code{cmake --build --preset debug-sanitize} and \code{ctest --preset debug-sanitize} reproduce exactly the same build
on every machine. Catch2's \code{catch\_discover\_tests} registers each \code{TEST\_CASE} as a separate CTest test.

% ------------------------------------------------------------------------------
\section{Testing}

\begin{keyidea}[a test must be able to fail]
For every test, ask: would it still pass if I flipped the sign of X, transposed Y, or swapped two blocks? If yes, the test is
decoration. Run it once against a deliberately broken implementation to make sure it fails.
\end{keyidea}

The project's test strategy layers several kinds of evidence.
\begin{enumerate}
  \item \textbf{Algebraic unit tests}: identities such as $\Log(\Exp(\phi)) = \phi$ at zero, tiny, generic and near-$\pi$ angles.
  \item \textbf{Finite-difference Jacobian tests} built on $\bplus/\bminus$, at non-trivial states, with tolerances justified in a comment.
  \item \textbf{Closed-form trajectory tests}: synthetic trajectories with analytic IMU data and convergence-order checks.
  \item \textbf{Statistical tests}: Monte Carlo covariance checks and NEES/NIS with $\chi^2$ bounds, seeded with
        \code{std::mt19937\_64}. Because \code{std::normal\_distribution} differs between standard libraries, never hard-code values from one
        run; test statistical properties.
  \item \textbf{Invariance tests}: cloning enabled versus disabled must give identical IMU estimates when no visual updates occur.
  \item \textbf{Resource tests}: allocation guards around every core hot path.
  \item \textbf{Sanitizers}: ASan and UBSan builds with \code{-fno-sanitize-recover=all}, so that any report fails the test run
        (\cref{sec:sanitizers}).
\end{enumerate}

% ------------------------------------------------------------------------------
\section{Debugging a diverging filter}

When the full system misbehaves, resist the urge to tune. Work down the layers instead.
\begin{enumerate}
  \item Run on \textbf{synthetic data} with exactly modeled noise. If it fails there, the bug is in the code.
  \item Disable visual updates. Does IMU-only propagation from ground truth behave as \cref{fig:cov-growth} predicts, with errors
        inside $3\sigma$?
  \item Replace visual updates by \textbf{pose updates} (Milestone 6). If those are consistent, the back end is sound.
  \item Feed the MSCKF \textbf{ground-truth feature tracks} from a synthetic scene. If that works, the problem is in the front end or the
        calibration.
  \item On real data, plot each residual component over time and look for bias, oscillation or outliers.
  \item Only then adjust noise parameters, and judge the result on consistency statistics.
\end{enumerate}

\section{Going further}

Once the baseline MSCKF runs, natural extensions are First-Estimates Jacobians (\cref{ch:observability}); online calibration of the
camera-IMU extrinsics and time offset~\cite{li2013msckf2, li2014timeoffset}; stereo observations from \code{cam1}; keeping long-lived
features as SLAM landmarks; zero-velocity updates; and evaluation on the harder EuRoC sequences (V1\_03, MH\_05) and other
datasets. Each extension adds states and Jacobians, and each Jacobian gets its derivation and its finite-difference test.

\begin{takeaways}
\begin{itemize}
  \item Heap allocation has unpredictable cost; the stack is cheap but small. Large workspaces are members, created once.
  \item Determinism requires zero allocation in the core: fixed-max-size matrices, member workspaces, proven by allocation guards and
        \code{EIGEN\_RUNTIME\_NO\_MALLOC} (Eigen allocates with \code{malloc}, which the guard does not see).
  \item Never use \code{auto} with Eigen expressions; use \code{.noalias()} only when there is no aliasing.
  \item Keep time in integer nanoseconds and process one presorted event stream.
  \item Sanitizers turn undefined behavior into reports; without \code{-fno-sanitize-recover=all} UBSan reports do not fail tests.
  \item Every test must be able to fail; layer algebraic, Jacobian, closed-form, statistical and resource tests.
  \item Debug by removing layers until the failure disappears, not by tuning.
\end{itemize}
\end{takeaways}

\begin{exercises}
  \item Write a minimal allocation guard by overriding the global \code{operator new} and \code{operator delete}, and show that it detects a
        \code{std::vector::push\_back} beyond capacity but not a fixed-max Eigen \code{resize}.
  \item Construct an example in which \code{auto} with an Eigen expression produces wrong numbers without crashing. Does
        AddressSanitizer report it?
  \item Measure the time of 10\,000 covariance propagation steps with fixed-max and with fully dynamic Eigen matrices, in both build
        presets. Explain the difference.
  \item Take one of your Jacobian tests, break the Jacobian in three different ways, and record which breakages the test detects.
  \item Write a function that reads one element past the end of a \code{std::array} and another that overflows an \code{int}. Build both presets and
        record what each sanitizer reports, and what the release build does.
  \item Pass a row-major block to a function taking \code{Eigen::Ref<const Eigen::MatrixXd>} inside an allocation guard. What happens, and what
        happens if the parameter is a writable \code{Ref}?
  \item Add a third alternative to the \code{SensorEvent} variant without handling it in the \code{Overloaded} visitor. Read the compiler error.
\end{exercises}
% Solutions to the exercises of Chapter 30.
\begin{solutions}
  \item The core of a minimal guard:
\begin{lstlisting}
inline int g_allocations = 0;                  // reset before, check after
void* operator new(std::size_t n) {
  ++g_allocations;
  if (void* p = std::malloc(n)) return p;
  throw std::bad_alloc();
}
void operator delete(void* p) noexcept { std::free(p); }
void operator delete(void* p, std::size_t) noexcept { std::free(p); }
\end{lstlisting}
        Reset the counter immediately before the code under test and read it immediately after, because the test framework allocates too. Measured
        with clang/libc++:
        \begin{itemize}
          \item four \code{push\_back} calls into a vector with \code{reserve(4)}: 0 allocations; the fifth: 1 (the reallocation);
          \item resizing a fixed-max $75\times75$ matrix from $15\times15$ to $75\times75$: 0, because the storage is inline;
          \item a dynamic \code{Eigen::MatrixXd(10, 10)}: also 0. This is not because it does not allocate, but because Eigen allocates with \code{std::malloc}
                through its own aligned allocator, which never calls \code{operator new}.
        \end{itemize}
        With \code{EIGEN\_RUNTIME\_NO\_MALLOC} and \code{set\_is\_malloc\_allowed(false)}, the same \code{MatrixXd} aborts with Eigen's assertion ``heap allocation is forbidden''. The guard and
        the flag are therefore both needed. Since the flag is defined only in the \code{debug-sanitize} preset, Eigen allocations are caught only there. Defining it for the
        test targets of both presets would close that gap; it is a design decision to record in \code{PROGRESS.md}. A complete guard also overrides the aligned
        forms \code{operator new(std::size\_t, std::align\_val\_t)}, which the library implements without calling the plain one. The C++ experiments of these
        solutions are in \code{scripts/solutions/cpp}, with \code{build\_and\_run.ps1} to build and run them in both configurations.

  \item Five candidates, each assigned afterwards to a concrete type and built in both presets:
        \begin{itemize}
          \item \code{auto e = R * (a + b);} gives correct numbers in both builds. The sum is a nested expression, which Eigen stores by value, so nothing dangles.
          \item \code{auto e = (a + b).eval().transpose();} prints the right numbers in release, by luck. ASan reports \code{stack-use-after-scope}.
          \item \code{auto e = make().head<2>();}, with \code{make()} returning a \code{Vector3d} by value, behaves the same way: right by luck in release, reported by ASan.
          \item \code{auto e = (X * 2.0).eval().block(0, 0, 2, 2);} with a dynamic \code{X} prints garbage ($10^{-311}$ entries) in release without crashing.
                ASan reports \code{stack-use-after-scope}.
          \item \code{auto e = A * v;} followed by a change to \code{A} produces the product with the \emph{new} \code{A} in both builds, and no sanitizer reports anything.
        \end{itemize}
        Dangling references are undefined behavior: the optimized build may print the right numbers because
        the dead temporary's stack slot has not been reused yet, or garbage, as in the dynamic case. AddressSanitizer reports them reliably, as a read of a variable
        whose scope has ended. The stale expression is not undefined behavior at all, since every access is to valid memory. The expression is simply evaluated
        later than the reader assumed, so no sanitizer can report it. Only the rule ``no \code{auto} with Eigen'' prevents it.

  \item Ten thousand propagation steps of a $75\times75$ covariance (the $15\times15$ block with $\Phi P\Phi\T + Q$, the $15\times60$ cross block with $\Phi P_{IC}$), per step:
        \begin{center}\small
        \begin{tabular}{@{}l c c@{}}
        \toprule
        implementation & release & debug-sanitize \\
        \midrule
        fixed-max, member workspaces & \SI{3.0}{\micro\second} & \SI{500}{\micro\second} \\
        dynamic, temporaries in expressions & \SI{3.3}{\micro\second} & \SI{525}{\micro\second} \\
        dynamic, preallocated workspaces & \SI{3.3}{\micro\second} & \SI{505}{\micro\second} \\
        \bottomrule
        \end{tabular}
        \end{center}
        In release, the arithmetic (about $4\times10^4$ flops) dominates. The three small heap allocations per step of the naive dynamic version cost a few hundred
        nanoseconds, and fixed sizes let Eigen unroll the $15\times15$ parts. That gives a 10\% difference, not an order of magnitude. The dynamic version with workspaces does not allocate
        either (it ran with Eigen allocations forbidden). The case for fixed-max matrices is therefore not average speed but the \emph{guarantee}: the storage is
        inline, its bound is checked, and no forgotten \code{.noalias()} or temporary can allocate in the hot path. That is what makes the worst case predictable. The
        sanitizer build is about 170 times slower. At \code{-O0} none of Eigen's small template functions is inlined, and ASan checks every load and store. Use it to
        prove correctness (both allocation-free versions ran with \code{malloc} forbidden), never to measure speed.

  \item With the right Jacobian $\Jr$ of \cref{ch:lie}, tested against central differences of $\Log(\Exp(\phi)\T\Exp(\phi + \delta))$ ($h = 10^{-6}$, so the correct version agrees to $10^{-10}$), the
        largest error of three broken versions at four test points is:
        \begin{center}\small
        \begin{tabular}{@{}l c c c c@{}}
        \toprule
        breakage & $\phi = 0$ & $\|\phi\| = 0.05$ & $\|\phi\| = 0.5$ & $\|\phi\| = 2.5$ \\
        \midrule
        sign of the $\sx{\phi}$ term flipped (this is $\Jl$) & 0 & $4\times10^{-2}$ & $0.4$ & $1.1$ \\
        second-order term dropped & 0 & $4\times10^{-4}$ & $4\times10^{-2}$ & $0.65$ \\
        Taylor branch used below $0.1$~rad & 0 & $4\times10^{-4}$ & $8\times10^{-11}$ & $9\times10^{-11}$ \\
        \bottomrule
        \end{tabular}
        \end{center}
        A test at $\phi = 0$ detects nothing, since every version equals $I$ there. A test at a single generic angle misses the wrong branch threshold. A loose tolerance
        ($10^{-3}$, say) misses the dropped term at small angles. Hence the rules: several random points spanning the regimes (tiny, inside every branch, generic, near
        $\pi$), and a tolerance derived from the finite-difference error rather than chosen for comfort.

  \item Reading \code{a[4]} of a \code{std::array<int, 4>} and computing \code{INT\_MAX + 1}:
        \begin{itemize}
          \item \emph{release}: the read returned 1634890337 (whatever lay next on the stack) and the addition printed $-2147483648$. No crash and no message.
                Neither result is guaranteed: both are undefined behavior, and the optimizer may, for example, fold \code{x + 1 > x} to \code{true};
          \item \emph{debug-sanitize}: ASan stops at the read with \code{stack-buffer-overflow ... READ of size 4} and the file and line of the access. UBSan stops at the
                addition with \code{runtime error: signed integer overflow: 2147483647 + 1 cannot be represented in type 'int'}. With \code{-fno-sanitize-recover=all}
                both abort, so a test that triggers them fails.
        \end{itemize}
        Note that \code{std::array::operator[]} does not check bounds. ASan caught the read because the array lives on the stack between redzones. Compiling with libc++'s
        hardening (\code{-D\_LIBCPP\_HARDENING\_MODE=\_LIBCPP\_HARDENING\_MODE\_FAST}) adds a cheap bounds check that traps even in an optimized build: the process stopped with an
        illegal-instruction exception.

  \item Binding a $3\times3$ block of a row-major matrix to \code{Eigen::Ref<const Eigen::MatrixXd>} compiles and gives the right result, but \code{Ref} cannot describe a row-major block
        with column-major strides, so it copies the block into an internal \code{MatrixXd}. The operator-new guard counts 0 allocations, because the copy is an Eigen
        allocation through \code{malloc}. With \code{set\_is\_malloc\_allowed(false)} the call aborts with Eigen's assertion. A column-major block binds without any copy (it passes with
        allocations forbidden). With a writable \code{Eigen::Ref<Eigen::MatrixXd>}, the call does not compile: \code{no matching function for call to 'scale'}. A writable reference cannot
        silently redirect writes into a copy. Keep one storage order throughout the core (column-major, Eigen's default), or write the \code{Ref} type for the layout
        actually passed.

  \item With \code{PoseEvent} added to the variant and no lambda for it, libc++ stops with
\begin{lstlisting}[language={}, basicstyle=\ttfamily\footnotesize]
error: static assertion failed due to requirement
  'is_invocable_v<Overloaded<...>, const PoseEvent &>':
  std::visit requires the visitor to be exhaustive.
\end{lstlisting}
        This is followed by notes leading back to the \code{std::visit} call. Every dispatch site that forgot the new event type fails to compile, which is the point of the
        design. The one way to defeat it is a catch-all generic lambda, \code{[](const auto\&) \{\}}, which silently swallows new alternatives. Do not put one in the
        event dispatch.
\end{solutions}

\fussy
